%! AP Statistics — Unit 9, Lesson 9.2 — Homework (Answer Key)
\documentclass[10pt]{article}
\usepackage{apstats-article}
\usepackage{apstats-key}

\begin{document}

\pageheader{Unit 9, Lesson 9.2}{Homework --- Answer Key}
\namedateperiod

\vspace{0.1in}

\textbf{Problem 1: Sleep and Reaction Time}

\vspace{0.05in}
A random sample of 22 college students was used to study the relationship between hours
of sleep per night ($x$) and reaction time in milliseconds ($y$). Computer output is
shown below.

\vspace{0.1in}
\begin{center}
\renewcommand{\arraystretch}{1.3}
\begin{tabular}{lrrrr}
    \textbf{Predictor} & \textbf{Coef} & \textbf{SE Coef} & \textbf{T} & \textbf{P} \\
    \hline
    Constant & 412.3 & 28.4 & 14.52 & 0.000 \\
    Sleep    & $-18.6$ & 5.3 & $-3.51$ & 0.002 \\
    \hline
    \multicolumn{5}{l}{$s = 24.1$ \quad $n = 22$}
\end{tabular}
\end{center}

\vspace{0.1in}
\begin{enumerate}[label=\alph*.]

    \item \ans{$b = -18.6$ ms per hour of sleep;\quad $SE_b = 5.3$.}

    \item \ans{\textbf{L}: Random scatter in residual plot --- linear condition met.
          \textbf{I}: Random sample of college students; $n = 22 \le 10\%$ of all students --- independent.
          \textbf{N}: $n = 22$; problem states conditions are met --- normal condition assumed.
          \textbf{E}: Roughly constant spread in residual plot --- equal variance met.
          \textbf{R}: Random sample --- random condition met.}

    \item \ans{$CI: -18.6 \pm 2.086(5.3) = -18.6 \pm 11.056 = (-29.656,\ -7.544)$ ms per hour.
          Round: $(-29.7,\ -7.5)$ ms per hour of sleep.}

    \item \ans{We are 95\% confident that the true slope of the population regression line
          relating reaction time to hours of sleep for college students is between $-29.7$
          and $-7.5$ milliseconds per hour of sleep.}

    \item \ans{Yes. The entire CI $(-29.7,\ -7.5)$ is below 0, so 0 is not in the interval.
          This provides evidence that $\beta < 0$: more hours of sleep is associated with
          faster (lower) reaction times --- a negative linear relationship.}

\end{enumerate}

\vspace{0.2in}

\textbf{Problem 2: Temperature and Heating Fuel Cost}

\vspace{0.05in}
A researcher studies 18 cities, recording average daily high temperature in January
($x$, ${}^\circ$F) and monthly heating fuel cost ($y$, \$/month).
Output gives $b = -4.8$, $SE_b = 1.1$, and $n = 18$.
Use $t^* = 2.120$ ($df = 16$). Assume all LINER conditions are met.

\vspace{0.1in}
\begin{enumerate}[label=\alph*.]

    \item \ans{$CI: -4.8 \pm 2.120(1.1) = -4.8 \pm 2.332 = (-7.132,\ -2.468)$ dollars per
          degree Fahrenheit. We are 95\% confident that the true slope of the population
          regression line relating monthly heating fuel cost to January temperature for
          all cities is between $-\$7.13$ and $-\$2.47$ per ${}^\circ$F.}

    \item \ans{Yes. The value $\beta = -5$ falls inside the CI $(-7.13,\ -2.47)$, so the
          utility company's claim is consistent with the data.}

    \item \ans{Narrower. The formula $SE_b = s/(s_x\sqrt{n-1})$ shows that a larger $n$
          produces a smaller $SE_b$. With $n = 72$ (four times as large), $SE_b$ would
          decrease, reducing the margin of error and producing a narrower CI.}

\end{enumerate}

\end{document}
